\documentclass{article}

% if you need to pass options to natbib, use, e.g.:
%     \PassOptionsToPackage{numbers, compress}{natbib}
% before loading neurips_2025

% The authors should use one of these tracks.
% Before accepting by the NeurIPS conference, select one of the options below.
% 0. "default" for submission
 % \usepackage{neurips_2025}
 \usepackage[main, final]{neurips_2025}
% the "default" option is equal to the "main" option, which is used for the Main Track with double-blind reviewing.
% 1. "main" option is used for the Main Track
%  \usepackage[main]{neurips_2025}
% 2. "position" option is used for the Position Paper Track
%  \usepackage[position]{neurips_2025}
% 3. "dandb" option is used for the Datasets & Benchmarks Track
 % \usepackage[dandb]{neurips_2025}
% 4. "creativeai" option is used for the Creative AI Track
%  \usepackage[creativeai]{neurips_2025}
% 5. "sglblindworkshop" option is used for the Workshop with single-blind reviewing
 % \usepackage[sglblindworkshop]{neurips_2025}
% 6. "dblblindworkshop" option is used for the Workshop with double-blind reviewing
%  \usepackage[dblblindworkshop]{neurips_2025}

% After being accepted, the authors should add "final" behind the track to compile a camera-ready version.
% 1. Main Track
 % \usepackage[main, final]{neurips_2025}
% 2. Position Paper Track
%  \usepackage[position, final]{neurips_2025}
% 3. Datasets & Benchmarks Track
 % \usepackage[dandb, final]{neurips_2025}
% 4. Creative AI Track
%  \usepackage[creativeai, final]{neurips_2025}
% 5. Workshop with single-blind reviewing
%  \usepackage[sglblindworkshop, final]{neurips_2025}
% 6. Workshop with double-blind reviewing
%  \usepackage[dblblindworkshop, final]{neurips_2025}
% Note. For the workshop paper template, both \title{} and \workshoptitle{} are required, with the former indicating the paper title shown in the title and the latter indicating the workshop title displayed in the footnote.
% For workshops (5., 6.), the authors should add the name of the workshop, "\workshoptitle" command is used to set the workshop title.
% \workshoptitle{WORKSHOP TITLE}

% "preprint" option is used for arXiv or other preprint submissions
 % \usepackage[preprint]{neurips_2025}

% to avoid loading the natbib package, add option nonatbib:
%    \usepackage[nonatbib]{neurips_2025}

\usepackage[utf8]{inputenc} % allow utf-8 input
\usepackage[T1]{fontenc}    % use 8-bit T1 fonts
\usepackage{hyperref}       % hyperlinks
\usepackage{url}            % simple URL typesetting
\usepackage{booktabs}       % professional-quality tables
\usepackage{amsfonts}       % blackboard math symbols
\usepackage{nicefrac}       % compact symbols for 1/2, etc.
\usepackage{microtype}      % microtypography
\usepackage{xcolor}         % colors
\PassOptionsToPackage{numbers, compress}{natbib}

% Note. For the workshop paper template, both \title{} and \workshoptitle{} are required, with the former indicating the paper title shown in the title and the latter indicating the workshop title displayed in the footnote. 
\title{JEPA‑NAV: A Vision-Based Joint Embedding Architecture for Safe Indoor Navigation}


% The \author macro works with any number of authors. There are two commands
% used to separate the names and addresses of multiple authors: \And and \AND.
%
% Using \And between authors leaves it to LaTeX to determine where to break the
% lines. Using \AND forces a line break at that point. So, if LaTeX puts 3 of 4
% authors names on the first line, and the last on the second line, try using
% \AND instead of \And before the third author name.


% \author{%
%   David S.~Hippocampus\thanks{Use footnote for providing further information
%     about author (webpage, alternative address)---\emph{not} for acknowledging
%     funding agencies.} \\
%   Department of Computer Science\\
%   Cranberry-Lemon University\\
%   Pittsburgh, PA 15213 \\
%   \texttt{hippo@cs.cranberry-lemon.edu} \\
%   % examples of more authors
%   % \And
%   % Coauthor \\
%   % Affiliation \\
%   % Address \\
%   % \texttt{email} \\
%   % \AND
%   % Coauthor \\
%   % Affiliation \\
%   % Address \\
%   % \texttt{email} \\
%   % \And
%   % Coauthor \\
%   % Affiliation \\
%   % Address \\
%   % \texttt{email} \\
%   % \And
%   % Coauthor \\
%   % Affiliation \\
%   % Address \\
%   % \texttt{email} \\
% }

\author{
Edward Ajayi\thanks{Carnegie Mellon University Africa, Kigali, Rwanda. \texttt{eaajayi@andrew.cmu.edu}} \\
CMU Africa, Kigali, Rwanda \\
\texttt{eaajayi@andrew.cmu.edu}
\And
Emrys Lasidlaus\thanks{Carnegie Mellon University Africa, Kigali, Rwanda. \texttt{eladisla@andrew.cmu.edu}} \\
CMU Africa, Kigali, Rwanda \\
\texttt{eladisla@andrew.cmu.edu}
}


\begin{document}


\maketitle


\begin{abstract}
  Navigating indoor environments is challenging for visually impaired individuals due to cluttered layouts, dynamic obstacles, and partial observability. We propose JEPA‑NAV, a vision-first system that leverages Joint Embedding Predictive Architectures (JEPA) to learn latent representations of indoor scenes. These embeddings support affordance prediction (walkable areas, obstacles, and risky zones) with associated uncertainty, which is integrated into a risk-aware navigation planner. By combining predictive latent vision with uncertainty estimation, JEPA‑NAV aims to improve navigation safety and generalization across previously unseen indoor layouts. Experiments will be conducted on publicly available 3D indoor datasets and simulators, establishing a baseline for future research in assistive, vision-based indoor navigation.
\end{abstract}


\section{Introduction}

Indoor navigation remains a significant challenge for visually impaired individuals seeking independent mobility in closed environments. Advances in vision-based technologies and vision-language models have spurred research on indoor navigation, yet effective systems must still perceive obstacles, identify walkable paths, and anticipate hazards while accounting for uncertainty. Navigation becomes particularly difficult when models encounter unfamiliar environments \cite{najjar2022dynamic}, as traditional approaches often fail to generalize and rarely provide safety-aware predictions.  

Self-supervised models, such as JEPA \cite{assran2025v}, have demonstrated the ability to learn predictive latent representations in vision tasks that generalize across diverse scenes. Motivated by these capabilities and their application in different vision tasks, we propose a vision-first, JEPA-based system that predicts environmental affordances along with associated uncertainty, enabling **safer and more reliable indoor navigation** for visually impaired users.


% \section{Related Works}
% Previous research has leveraged vision capabilities in providing navigation support for visually impaired using different techniques. 

% ~\cite{li2018vision} proposed a vision-based indoor navigation system for visually impaired users. It builds a semantic indoor map, performs real-time obstacle detection using an RGB-D camera, and provides multimodal feedback via audio and haptics. Field tests showed improved navigation safety and independence\cite{li2018vision}. 

\section{Related Works}

Previous research has explored vision-based approaches to support indoor navigation for visually impaired users. ~\cite{li2018vision} proposed a system that constructs semantic indoor maps, detects obstacles in real time using an RGB-D camera, and provides multimodal feedback through audio and haptics, demonstrating improved navigation safety and independence. ~\cite{na2024improving} focused on walking path generation for a robotic SmartCane, integrating human motion constraints and motion primitives to guide users efficiently and naturally in indoor environments. Both studies highlight the importance of real-time perception and path planning for safe assistive navigation, though they rely heavily on specialized hardware and system-specific configurations.

Recent advances in predictive visual representations provide new opportunities for generalizable and anticipatory navigation. ~\cite{miao_jepa-vla_2026} introduced JEPA-VLA, integrating predictive video embeddings into vision-language-action models to capture task-relevant temporal dynamics while ignoring unpredictable environmental factors, improving sample efficiency and generalization. Similarly, ~\cite{assran2025v} proposed V-JEPA 2, a self-supervised video model trained on large-scale video data to learn joint predictive embeddings capable of understanding, predicting, and planning in physical environments. These predictive latent representations offer a promising foundation for vision-first indoor navigation systems, as they can anticipate obstacles and affordances in previously unseen or dynamic environments, forming the basis for our proposed JEPA‑NAV approach.

\section{Proposed Methodology}

Our proposed approach, \textbf{JEPA‑NAV}, is a vision-first framework for safe indoor navigation for visually impaired users. The methodology consists of the following key components:

\begin{enumerate}
    \item \textbf{Visual Representation Learning:} We leverage JEPA-based self-supervised embeddings to learn predictive latent representations from large-scale video and image datasets. These embeddings capture environment structure, dynamics, and temporal patterns, providing a foundation for anticipating obstacles and walkable areas.

    \item \textbf{Affordance Prediction:} Using the learned visual embeddings, we predict environmental affordances such as walkable zones, obstacles, and potential hazards. This module produces a spatial map of navigable paths, allowing the system to reason about which areas are safe for movement.

    \item \textbf{Uncertainty Modeling:} To account for unfamiliar or dynamic environments, we integrate uncertainty estimation into the affordance predictions. Probabilistic outputs or ensemble approaches allow the system to quantify confidence in its perception, guiding risk-aware path planning.

    \item \textbf{Navigation and Path Planning:} The predicted affordances and uncertainty maps are used to generate safe paths for indoor navigation. Paths are selected to minimize collision risk and maximize navigable area, enabling the user to traverse unfamiliar environments with greater safety.

    \item \textbf{Optional Multimodal Integration:} While the focus is vision-first, additional inputs such as depth or audio cues can be incorporated to enhance robustness, providing a multimodal reasoning layer for future extensions.
\end{enumerate}

This methodology emphasizes our aim of generalization to unseen environments, predictive understanding of obstacles, and safety-aware navigation, forming the basis for experimental evaluation in simulation and, potentially, low-cost real-world devices.

\section{Datasets and Experimental Setup}

To evaluate the proposed JEPA‑NAV system, we plan to use the following datasets and simulation environments, which provide realistic indoor scenes, visual information, and affordance labels where available:

\begin{itemize}
    \item \textbf{Matterport3D} \cite{yadav2023habitat}: A large-scale dataset of indoor environments with RGB-D images and 3D reconstructions. It provides panoramic images, depth maps, and semantic annotations, suitable for training visual representations and predicting walkable areas.

    \item \textbf{Replica} \cite{straub2019replicadatasetdigitalreplica}: High-quality, photorealistic 3D indoor reconstructions with semantic labels. Replica is particularly useful for evaluating affordance prediction and navigation in environments with varied layouts and obstacles.

    \item \textbf{Gibson} \cite{yokoyama2022benchmarking}: A dataset and simulator of real-world indoor scenes with navigation meshes, RGB-D data, and realistic physics. It allows testing navigation pipelines under dynamic conditions and measuring collision rates, path efficiency, and coverage.

    % \item \textbf{MINOS / Habitat Simulator} \cite{habitat}: Simulation environments that integrate the above datasets and allow **embodied agent testing**. They provide the ability to test safety-aware path planning, obstacle avoidance, and generalization to unseen indoor environments.

    \item \textbf{Affordance-related datasets} (optional for training affordance prediction):
    \begin{itemize}
        \item \textbf{ADE20K} \cite{zhou2017ade20k}: Provides scene parsing and object segmentation labels that can be mapped to walkable vs non-walkable regions.  
        \item \textbf{MP3D Affordances / RoboTHOR variants} \cite{deitke2020robothor}: Some extensions provide annotations for traversable regions, interactable objects, and hazard zones.
    \end{itemize}
\end{itemize}

\paragraph{Experimental Setup}  
The system will be trained and evaluated in simulation using Habitat or MINOS, leveraging GPU resources for JEPA embedding pretraining and affordance prediction modules. Metrics include:
\begin{itemize}
    \item \textbf{Navigation success rate:} Percentage of successfully reached goals without collisions.  
    \item \textbf{Collision rate:} Number of collisions per trajectory.  
    \item \textbf{Affordance prediction accuracy:} Accuracy in predicting walkable areas and obstacles.  
    \item \textbf{Uncertainty calibration:} Measuring how well predicted confidence aligns with actual safety in navigation.  
\end{itemize}

This setup allows systematic evaluation of generalization, predictive understanding, and safety-aware navigation without requiring extensive real-world data collection.

\section{Expected Results / Evaluation}

We plan to evaluate JEPA‑NAV on both simulation and benchmark datasets to demonstrate its effectiveness in safe indoor navigation for visually impaired users. Key evaluation objectives include:

\begin{itemize}
    \item \textbf{Safe navigation in unseen environments:} Measuring success rate in reaching goals without collisions in previously unobserved indoor layouts.
    \item \textbf{Accurate affordance prediction:} Evaluating the system's ability to correctly identify walkable paths, obstacles, and potential hazards.
    \item \textbf{Uncertainty calibration:} Assessing how well predicted confidence aligns with actual safety, guiding risk-aware path planning.
    \item \textbf{Optional user simulation:} Testing navigation performance using blindfolded agents or virtual users to approximate real-world usage scenarios.
\end{itemize}


\section{Expected Contributions}

This work aims to make the following contributions:

\begin{itemize}
    \item \textbf{First application of JEPA-based predictive embeddings for assistive indoor navigation:} enabling anticipatory understanding of the environment and supporting predictive decision-making.
    \item \textbf{Uncertainty-aware affordance prediction module:} informing risk-aware path planning to improve navigation safety for visually impaired users.
    \item \textbf{Vision-first navigation pipeline:} capable of generalizing to previously unseen indoor environments without task-specific retraining.
    \item \textbf{Establishment of a reproducible baseline:} providing a foundation for future research in predictive vision systems for assistive navigation.
\end{itemize}

\section{Conclusion / Significance}

JEPA‑NAV aims to advance indoor navigation systems for visually impaired users by combining predictive, self-supervised visual embeddings with uncertainty-aware path planning. The significance of this work includes:

\begin{itemize}
    \item Enhancing independence and safety for visually impaired individuals in indoor spaces.
    \item Leveraging JEPA-based predictive embeddings and uncertainty modeling to anticipate obstacles and affordances in unfamiliar environments.
\end{itemize}

\section{Timeline / Work Plan}

The project is planned over a 10-week period as follows:

\begin{table}[htbp]
\centering
\begin{tabular}{|c|l|}
\hline
\textbf{Week(s)} & \textbf{Tasks} \\ \hline
1--2 & Literature review: study related work on JEPA, vision-based navigation, and affordance prediction. \\ \hline
3--4 & Integrating V-JEPA-2 embeddings into the pipeline \& evaluating representation quality for indoor scenes \\ \hline
5--6 & Developing the affordance prediction and uncertainty modeling modules. \\ \hline
7--8 & Integrating modules into the navigation pipeline and testing in simulation (Habitat/MINOS). \\ \hline
9 & Conducting evaluation: navigation performance, affordance accuracy, and uncertainty calibration. \\ \hline
10 & Analysis of results, refinement of pipeline, and preparation of final write-up for submission. \\ \hline
\end{tabular}
\caption{10-Week Work Plan for JEPA‑NAV Project}
\label{tab:workplan}
\end{table}

\bibliographystyle{plainnat} % Or another style like 'abbrvnat'
\bibliography{references}    % Assuming your file is references.bib



\end{document}